\chapter{Teorema de confluencia HMT–MD: forma electrogravitatoria, sistema Einstein–Schrödinger y coordinación de acción, materia, campos, geometría e información}
\label{chap:parte-iv-51}

La unificación se expresa como compatibilidad vertical de publicaciones distintas de un mismo estado: acoplamiento y vacío, acción y gravedad, geometría e información, pantalla y torsión, reloj y horizonte. El teorema conserva los dominios y mapas de cada respuesta.

\section{Teorema de confluencia de las respuestas dimensionales}
\label{rm:ch:confluencia-general}

\subsection{Antecedente com\'un y producto fibrado}

Los resultados del volumen proceden de lectores heterog\'eneos aplicados a
un mismo estado geneal\'ogico enriquecido.  Sea

\begin{equation}
 X=\bigl(X_{\APP},X_{\TRIT},X_{\TPK},
          \mathcal C,\mathcal B,\mathcal R,\mathcal M,\partial X\bigr)
 \label{rm:eq:confluence-state}
\end{equation}

el estado que conserva cilindros, acarreo, ruta, memoria y frontera. Sobre
\(X\) se construyen los lectores tipados

\begin{equation}
 P_i:X\longrightarrow Y_i,
 \qquad
 i\in I:=\{\alpha,{\rm act},{\rm ang},108,\square,{\rm obs},{\rm Per}\}.
 \label{rm:eq:confluence-objects}
\end{equation}

Sus codominios representan, respectivamente, el generador dodecaf\'asico de
acoplamiento, la pareja de acci\'on, el operador angular bicapa, el reloj
\(108\), la pantalla c\'ubica, el observador interno y el modo de Perron.
Para cada lector se toma su grafo, provisto de la proyecci\'on can\'onica

\begin{equation}
 \Gamma_i:=\{(x,y_i)\in X\times Y_i:P_i(x)=y_i\},
 \qquad
 \pi_i:\Gamma_i\longrightarrow X,
 \quad (x,y_i)\longmapsto x.
 \label{rm:eq:confluence-reader-graphs}
\end{equation}

El objeto maestro es el producto fibrado bien tipado de esos grafos:

\begin{equation}
 \boxed{
 \mathfrak M_X=
 \mathop{\times}_{X,\,i\in I}\Gamma_i
 =\left\{\bigl(x,(y_i)_{i\in I}\bigr):
 P_i(x)=y_i\ \text{para todo }i\in I\right\}.}
 \label{rm:eq:master-fiber-product}
\end{equation}

Las proyecciones \(\pi_i\) fijan materialmente el mismo antecedente. Un
elemento de \(\mathfrak M_X\) es una familia de salidas compatibles y
acompa\~nadas por el estado que las genera.

\begin{definition}[Publicaci\'on dimensional]
Una publicaci\'on dimensional es una composici\'on
\begin{equation}
 \mathfrak M_X\xrightarrow{\;P\;}\mathcal L
 \xrightarrow{\;\operatorname{ev}_{\mathcal U}\;}\RR,
 \label{rm:eq:dimensional-publication}
\end{equation}
donde \(\mathcal L\) es una recta dimensional tipada y
\(\operatorname{ev}_{\mathcal U}\) es la evaluaci\'on final en una carta de
unidades. El mapa \(P\) construye primero el invariante y su genealog\'ia.
\end{definition}

\subsection{Independencia horizontal y compatibilidad vertical}

Dos lectores hermanos comparten antecedente y pueden satisfacer relaciones
cruzadas sin factorizar uno por el otro.  Esta distinci\'on gobierna las
dobles hojas del vac\'io, las realizaciones positiva y orientada, la doble
proyecci\'on arquimediana y excepcional, y las lecturas geom\'etrica,
informacional y de sabor del tr\'ansito \(6\to8\).

\begin{theorem}[Confluencia sin factorizaci\'on horizontal]
Sean \(P_i:\mathfrak M_X\to Y_i\), \(i=1,\ldots,r\), publicaciones
construidas sobre el mismo estado. Sup\'ongase que cada \(P_i\) conserva sus
datos tipados y que las relaciones \(R_j(P_1,\ldots,P_r)=0\) se demuestran
despu\'es de construir los lectores.  Entonces el morfismo conjunto
\begin{equation}
 P=(P_1,\ldots,P_r):\mathfrak M_X\longrightarrow
 Y_1\times\cdots\times Y_r
 \label{rm:eq:joint-publication}
\end{equation}
es compatible verticalmente con \(X\). Una factorizaci\'on
\(P_i=F_{ij}\circ P_j\) requiere, adem\'as, que \(P_j\) separe todas las
fibras que \(P_i\) distingue. La mera existencia de \(R_j=0\) no implica
esa separaci\'on.
\end{theorem}

\begin{proof}
La compatibilidad vertical se obtiene por la propiedad universal del
producto fibrado de los grafos \(\Gamma_i\): todas las componentes poseen
la misma imagen en \(X\).
Si \(P_i=F_{ij}\circ P_j\), entonces
\(P_j(x)=P_j(x')\) fuerza \(P_i(x)=P_i(x')\). Por tanto, cualquier par
\(x,x'\) que comparta la lectura \(j\) y difiera en la lectura \(i\) refuta
la factorizaci\'on. Una relaci\'on escalar \(R_j=0\) s\'olo restringe la imagen
conjunta; no garantiza la inclusi\'on de n\'ucleos necesaria para factorizar.
\end{proof}

El teorema explica por qu\'e el producto de hojas puede fijar el cono causal
mientras el cociente conserva orientaci\'on, y por qu\'e el \`area total puede
coexistir con un Hamiltoniano modular que distingue la procedencia \(90/120\).

\subsection{1.ª confluencia: acoplamiento, vacío y sector electrodébil}

Sea

\begin{equation}
 D=D_A=\det T,
 \qquad
 Q(D)=-\frac9{25}\log D=4\pi\alpha_{\HMT}.
 \label{rm:eq:confluence-q-d}
\end{equation}

La memoria orientada se expresa mediante

\begin{equation}
 \rho=e^{2C^*_{\rm rad}},
 \qquad
 q_-=\sqrt{D\rho},
 \qquad
 q_+=\sqrt{D/\rho}.
 \label{rm:eq:confluence-qpm}
\end{equation}

El c\'alculo funcional

\begin{equation}
 F(z)=\frac{1-z^{90}}{1-z^{120}},
 \qquad
 r_\pm=F(q_\pm)
 \label{rm:eq:confluence-F}
\end{equation}

produce

\begin{equation}
 \widehat\mu=r_-^2,
 \quad
 \widehat\varepsilon=r_+^2,
 \quad
 \widehat Z=\frac{r_-}{r_+},
 \quad
 \widehat c=(r_-r_+)^{-1}.
 \label{rm:eq:confluence-vacuum}
\end{equation}

En la interfaz electrod\'ebil de \`arbol,

\begin{equation}
 g^2=\frac{Q(D)}{1-D},
 \qquad
 g'^2=\frac{Q(D)}{D},
 \label{rm:eq:confluence-gauge-couplings}
\end{equation}

de donde

\begin{equation}
 \boxed{
 (1-D)g^2=Dg'^2=Q(D),
 \qquad
 \frac1{g^2}+\frac1{g'^2}=\frac1{Q(D)}.}
 \label{rm:eq:confluence-ew-master}
\end{equation}

La primera igualdad publica el mismo acoplamiento en dos orientaciones de
calibre; la segunda lo recupera como suma de compliancias. El determinante
\(D\) gobierna la mezcla par y \(\rho\) conserva la orientaci\'on impar.

La coordenada de Higgs incorpora esta \`ultima:

\begin{equation}
 \lambda_H=\frac{Q(D)}{8(1-D)}D^{-3}\rho^{5/3}.
 \label{rm:eq:confluence-higgs}
\end{equation}

Definiendo

\begin{equation}
 \Xi=\frac{8(1-D)D^3\lambda_H}{Q(D)},
 \label{rm:eq:confluence-xi}
\end{equation}

se reconstruye

\begin{equation}
 \rho=\Xi^{3/5},
 \qquad
 \widehat Z=
 \frac{F(\sqrt D\,\Xi^{3/10})}
      {F(\sqrt D\,\Xi^{-3/10})}.
 \label{rm:eq:confluence-vacuum-higgs}
\end{equation}

La ecuaci\'on \cref{rm:eq:confluence-vacuum-higgs} constituye un control cruzado
entre orientaci\'on escalar e impedancia del vac\'io.  Su fuerza procede de la
construcci\'on paralela de ambos lectores sobre \(X\).

\subsection{2.ª confluencia: acción, gravedad y forma electrogravitatoria}

Sea

\begin{equation}
 Q_P=E_P,
 \qquad
 \ell_P,
 \qquad
 Q_P\ell_P=\hbar c,
 \qquad
 \frac{\ell_P}{Q_P}=\frac{G}{c^4}.
 \label{rm:eq:confluence-planck-pair}
\end{equation}

El producto es el cuanto de intercambio \(\hbar c\); el cociente es la
compliancia gravitatoria.  Introducimos

\begin{equation}
 \mathscr T_P=\frac{Q_P}{\ell_P}=\frac{c^4}{G},
 \qquad
 \mathscr C_P=\frac{\ell_P}{Q_P}=\frac{G}{c^4}.
 \label{rm:eq:confluence-tension-compliance}
\end{equation}

La ecuaci\'on de Einstein adopta la forma constitutiva

\begin{equation}
 G_{\mu\nu}+\Lambda g_{\mu\nu}
 =8\pi\mathscr C_P T_{\mu\nu}.
 \label{rm:eq:confluence-einstein-constitutive}
\end{equation}

Para una carga \(q\), la coordenada el\'ectrica elevada al plano de energ\'ia
es

\begin{equation}
 \mathcal Q(q)=Q_P\sqrt{\frac{Z}{4\pi\hbar}}\,q.
 \label{rm:eq:confluence-electric-coordinate}
\end{equation}

Entonces Newton y Coulomb se re\'unen en la forma bilineal

\begin{equation}
 \boxed{
 V(r)=\frac{\mathscr C_P}{r}
 \bigl(\mathcal Q_1\mathcal Q_2-E_1E_2\bigr).}
 \label{rm:eq:confluence-newton-coulomb}
\end{equation}

El cono \(\mathcal Q^2=E^2\) reproduce la condici\'on est\'atica de
extremalidad en la carta Einstein--Maxwell.  El producto de hojas del vac\'io
gobierna el cono causal; su cociente gobierna la elevaci\'on el\'ectrica y la
orientaci\'on.

La identidad cruzada con el sector electrod\'ebil es

\begin{equation}
 \boxed{
 \frac{\hbar^2}{8\pi\ell_P^2}
 \frac{(\kappa/\mu)}{e^2}
 =\frac1{4\pi\alpha_{\HMT}}
 =\frac1{g^2}+\frac1{g'^2}.}
 \label{rm:eq:confluence-einstein-maxwell-ew}
\end{equation}

Esta igualdad conserva su valor bajo la inversi\'on de la hoja de acci\'on:
\(G\) y \(e^2\) transforman contragredientemente, mientras
\(G\hbar\), \(Ge^2\) y \(\ell_P^2\) permanecen invariantes.

\subsection{3.ª confluencia: geometría, información y mezcla}

Sea \(N=N_{90}+N_{120}\) el n\'umero total de excitaciones de borde y

\begin{equation}
 D_{90/120}=90N_{90}+120N_{120}.
 \label{rm:eq:confluence-modular-provenance}
\end{equation}

El \`area lee \(N\), mientras el Hamiltoniano hologr\'afico modular lee
\(D_{90/120}\). Como

\begin{equation}
 \det\begin{pmatrix}1&1\\90&120\end{pmatrix}=30,
 \label{rm:eq:confluence-modular-determinant}
\end{equation}

la pareja reconstruye exactamente

\begin{equation}
 N_{90}=4N-\frac{D_{90/120}}{30},
 \qquad
 N_{120}=\frac{D_{90/120}}{30}-3N.
 \label{rm:eq:confluence-sector-reconstruction}
\end{equation}

La geometr\'ia total y la energ\'ia modular ponderada conservan, por tanto,
cantidad y procedencia.

El estado bicapa angular

\begin{equation}
 \rho_y=\frac{e^{-yR}}{2\cosh y}
 \label{rm:eq:confluence-rhoy}
\end{equation}

posee entrop\'ia par y divergencia orientada

\begin{equation}
 S(\rho_y)=\log(2\cosh y)-y\tanh y,
 \qquad
 D(\rho_y\Vert\rho_{-y})=2y\tanh y.
 \label{rm:eq:confluence-info-orientation}
\end{equation}

La involuci\'on \(y\mapsto-y\) se representa en el sector de sabor mediante

\begin{equation}
 R_{\rm CKM}
 =M_{\rm CKM}\operatorname{diag}(1,-1,1)M_{\rm CKM}^{-1},
 \qquad
 R_{\rm CKM}^2=I,
 \qquad
 \det R_{\rm CKM}=-1.
 \label{rm:eq:confluence-rckm}
\end{equation}

Barbero lee la componente impar de la transici\'on \(6\to8\), la entrop\'ia
lee su distribuci\'on par y \(R_{\rm CKM}\) realiza la misma inversi\'on en el
codominio de mezcla.  Son representaciones coordinadas de una bisagra com\'un,
con dominios y observables propios.

\subsection{4.ª confluencia: pantalla, torsión y memoria helicoidal}

En el cociente lorentziano \(Q_{\square}\), la soldadura y la respuesta
positiva satisfacen

\begin{equation}
 \boxed{
 (\beta_m^{\square})^*\Lambda_m^{\square}\beta_m^{\square}
 =\frac{56}{81}I_{Q_{\square}}.}
 \label{rm:eq:confluence-screen-energy}
\end{equation}

En el registro profundo, el operador de tornillo verifica

\begin{equation}
 \boxed{
 V^{-1}\mathscr SV=(2+\sqrt3)\ii\,\mathscr S.}
 \label{rm:eq:confluence-screw}
\end{equation}

La primera identidad mide el coste m\'inimo de respuesta de la pantalla; la
segunda transporta orientaci\'on y profundidad. Catalan aparece como el
segundo momento del componente cuaternario. Si
\(N_{\rm dep}\delta_n=n\delta_n\) denota el operador de profundidad del
registro helicoidal, entonces

\begin{equation}
 G_{\rm Cat}=\Tr(N_{\rm dep}^{-2}\mathcal Q_{\chi,-4}).
 \label{rm:eq:confluence-catalan}
\end{equation}

La ley de Perron

\begin{equation}
 M_n=\left(\frac{a_n}{a_0}\right)^{-6}
 \label{rm:eq:confluence-sixth-law}
\end{equation}

ofrece la covariancia exacta de una densidad cuadr\'atica de esp\'in.  El
realizador de Cartan debe transportar
\cref{rm:eq:confluence-screen-energy,rm:eq:confluence-screw} a una corriente y una
contorsi\'on cuyo modo homog\'eneo satisfaga
\cref{rm:eq:confluence-sixth-law}.  Este encadenamiento sit\'ua la interfaz
f\'isica en una flecha \`unica y falsable.

\subsection{Confluencia entre periodicidad del reloj, energía de horizonte y coste informacional}

El reloj \(108\) fija

\begin{equation}
 E_P=k_B\Theta_{\rm clk}\log3.
 \label{rm:eq:confluence-planck-trit}
\end{equation}

Para el parche de Sitter de radio \(R=\ell_P\), la energ\'ia de vac\'io, la
temperatura y la entrop\'ia cumplen

\begin{equation}
 E_\Lambda=T_H S_H=\frac{E_P}{2}
 =\frac{k_B\Theta_{\rm clk}\log3}{2}.
 \label{rm:eq:confluence-horizon-half}
\end{equation}

De aqu\'i resultan

\begin{equation}
 E_\Lambda t_P=\frac{\hbar}{2},
 \qquad
 E_\Lambda T_{108}=\frac{h}{2},
 \qquad
 \frac{S_H}{k_B}=\pi.
 \label{rm:eq:confluence-half-action}
\end{equation}

El volumen del vac\'io, el \`area del horizonte, la temperatura del ciclo y el
coste informacional de un trit publican la misma energ\'ia.  La multiplicidad
\(e^{S_H/k_B}=e^\pi\) coincide exactamente con el autovalor expansivo de
\(e^{\pi J_\varphi}\); la promoci\'on de esta igualdad a equivalencia de
microestados exige un entrelazador que conserve medida y din\'amica.

\subsection{Confluencia fibrada de las respuestas HMT–MD: vacío, acción, gravedad, información, mezcla, memoria y cosmología}

\begin{theorem}[Confluencia de respuestas HMT--MD]
Fijado el estado enriquecido \(X\) y las primitivas declaradas en cada lector,
las publicaciones de vac\'io, acci\'on, gravedad, informaci\'on, mezcla,
pantalla, memoria helicoidal y cosmolog\'ia son las proyecciones o
composiciones definidas sobre el producto fibrado de grafos
\cref{rm:eq:master-fiber-product}. Sus relaciones cruzadas quedan
generadas por
\cref{rm:eq:confluence-ew-master,rm:eq:confluence-newton-coulomb,%
rm:eq:confluence-einstein-maxwell-ew,rm:eq:confluence-sector-reconstruction,%
rm:eq:confluence-info-orientation,rm:eq:confluence-screen-energy,%
rm:eq:confluence-screw,rm:eq:confluence-horizon-half}.
Cada evaluaci\'on num\'erica es terminal y cada memoria orientada permanece en
el factor que la genera.
\end{theorem}

\begin{proof}
La construcci\'on de \(\mathfrak M_X\) fija el antecedente com\'un. Las
ecuaciones de vac\'io siguen del c\'alculo funcional de \(T\); las ecuaciones
electrod\'ebiles siguen de \(D\) y \(Q(D)\); la forma electrogravitatoria sigue
de las dos identidades de la pareja Planck; la reconstrucci\'on \(90/120\)
resulta de invertir una matriz entera de determinante \(30\); la orientaci\'on
informacional sigue del logaritmo de \(\rho_y\); la energ\'ia de pantalla
resulta de la respuesta \`unica en \(Q_{\square}\); el tornillo sigue de las
dos covariancias de Weyl y Pell; y la igualdad de horizonte sigue de las
f\'ormulas geom\'etricas de Sitter especializadas en \(R=\ell_P\). Ninguna
prueba utiliza el decimal terminal para seleccionar la rama.  El teorema de
independencia horizontal impide convertir las relaciones cruzadas en
factorizaciones no demostradas.
\end{proof}

\subsection{Corolario de confluencia metrológica entre las realizaciones HMT y SI}

Toda salida admite la qu\'intupla inseparable

\begin{equation}
 \mathcal K(P)=
 \bigl(
 \text{antecedente},
 \text{ecuaci\'on},
 \text{invariante},
 \text{significado},
 \text{valor terminal}
 \bigr).
 \label{rm:eq:confluence-quintuple}
\end{equation}

El atlas metrol\'ogico del volumen enumera estas qu\'intuplas.  Su contenido
probatorio es el siguiente: el antecedente fija la genealog\'ia; la ecuaci\'on
fija la operaci\'on; el invariante demuestra qu\'e sobrevive; el significado
explicita la lectura f\'isica; el valor terminal permite el reconocimiento
externo.  Separar una columna destruye informaci\'on causal; conservar las
cinco permite insertar el resultado en una carta de unidades sin convertir
esa carta en fundamento.

\begin{proofstatusbox}
El producto fibrado, las identidades algebraicas y el teorema de independencia
horizontal constituyen una formalizaci\'on exacta.  Las ecuaciones internas
citadas conservan el estatuto demostrado en los cap\'itulos que contienen sus demostraciones completas.  Las
lecturas Einstein--Maxwell, de Sitter, Einstein--Cartan y electrod\'ebil se
mantienen en sus interfaces declaradas cuando requieren un realizador
adicional.  El teorema organiza esas capas sin promoverlas por similitud
num\'erica.
\end{proofstatusbox}

\begin{falsifierbox}
La confluencia se refuta si dos componentes atribuidas al mismo elemento de
\(\mathfrak M_X\) proceden realmente de estados incompatibles, si falla una
de las identidades generadoras, si una evaluaci\'on terminal interviene en la
selecci\'on de su propio antecedente, o si se pierde una memoria que otro
lector distingue.  Una pretendida factorizaci\'on horizontal queda refutada
por cualquier par de estados con igual lector intermedio y salida final
distinta.  Las promociones f\'isicas se contrastan, adem\'as, con sus propios
realizadores, identidades de gauge, Bianchi, refinamiento y condiciones de
frontera.
\end{falsifierbox}


\section{Realización física discreta mediante producto fibrado y entrelazamiento de Cartan}

El producto fibrado común, sus componentes positivas y orientadas, la
localización reversible y el criterio de promoción suave proporcionan el
entrelazador material entre la confluencia abstracta y sus realizaciones
físicas. El desarrollo se conserva íntegro y con sus obstrucciones explícitas.

\begingroup
\let\chapter\section
\let\section\subsection
\let\subsection\subsubsection
\let\subsubsection\paragraph
\let\clearpage\relax
\let\cleardoublepage\relax
\section[Realización física discreta sobre rutas TPK]{Realización física
discreta: lecturas positiva y orientada sobre rutas TPK, dominio común y
localización reversible}
\label{sec:int68-realizacion-fisica-discreta}

La forma espectral, la realización de Hilbert, la sección dimensional y el
transporte de Cartan no son decoraciones añadidas a una salida clásica. Son
cuatro representaciones coordinadas de una misma ruta TPK enriquecida. Para
expresar esa unidad causal debe construirse primero su dominio común; no
basta escribir cuatro flechas con el mismo símbolo de partida.

\subsection{Dominio común como producto fibrado}

Sea \(\mathsf P_{\rm TPK}^{\rm enr}\) la categoría de estados TPK preparados
y rutas registradas compatibles con sus extremos, frontera y restricciones.
Sean
\[
 \mathsf D_{\rm sp}:=\mathsf{SurvPath}_{\rm HMT},\qquad
 \mathsf D_Q:=\operatorname{Dom}(\mathcal Q),\qquad
 \mathsf D_D:=\operatorname{Dom}(\mathfrak D_{\rm HMT}),\qquad
 \mathsf D_C:=\operatorname{Dom}(\rho_C^{\rm disc}),
\]
con sus funtores de olvido fieles
\[
 U_i:\mathsf D_i\longrightarrow\mathsf P_{\rm TPK}^{\rm enr}
 \qquad(i\in\{{\rm sp},Q,D,C\}).
\]
Cada \(U_i\) olvida únicamente la realización y conserva el estado y la ruta
TPK subyacentes. Definimos
\begin{equation}
 \mathsf P_{\rm CT}^{+}
 :=
 \mathsf D_{\rm sp}
 \times_{\mathsf P_{\rm TPK}^{\rm enr}}\mathsf D_Q
 \times_{\mathsf P_{\rm TPK}^{\rm enr}}\mathsf D_D
 \times_{\mathsf P_{\rm TPK}^{\rm enr}}\mathsf D_C.
 \label{eq:int68-pct-pullback}
\end{equation}
Un objeto es, por tanto, un estado preparado provisto de las cuatro
realizaciones sobre el mismo antecedente; una flecha es una ruta registrada
cuya imagen subyacente coincide en los cuatro factores. Si las
presentaciones se identifican sólo por isomorfismo, se usa el producto
fibrado bicategórico correspondiente; en la presentación fijada del registro
se adopta el producto fibrado estricto de
\eqref{eq:int68-pct-pullback}.

\begin{proposition}[Existencia de un objeto en el dominio común]
\label{prop:int68-pct-no-vacio}
La categoría \(\mathsf P_{\rm CT}^{+}\) no es vacía. No es necesario
particularizar la ruta electrónica central para probarlo: cualquier sección
coinductiva certificada, equipada con su flecha identidad y sus datos neutros, admite
levantamientos compatibles a los cuatro dominios de
\eqref{eq:int68-pct-pullback}.
\end{proposition}

\begin{proof}
El \cref{thm:generacion-monodromica-conjunta-rev7} proporciona, por ejemplo,
una sección superviviente preparada \(x_\pi\). Su identidad
\(\operatorname{id}_{x_\pi}\) es simultáneamente: una ruta superviviente;
un morfismo de \(\mathsf P_{\rm TPK}\), que \(\mathcal Q\) envía a la
identidad acotada de \(\mathcal Q(x_\pi)\); la ruta vacía con firma nula,
que el \cref{thm:funtor-dimensional-minimo} envía a
\((0,0,0,1)\); y la identidad del grupoide, que el
\cref{thm:realizacion-discreta-cartan-rev6} envía al elemento neutro
\((I,0)\) de Poincaré. Estos cuatro levantamientos olvidan al mismo objeto
\(x_\pi\) y al mismo morfismo \(\operatorname{id}_{x_\pi}\). Su cuádrupla
define, por la propiedad del producto fibrado, un objeto de
\(\mathsf P_{\rm CT}^{+}\).
\end{proof}

\begin{definition}[Adaptadores al dominio común]
\label{def:int68-pct-positive}
Los cuatro adaptadores son las proyecciones canónicas
\[
 \begin{aligned}
 F_{\rm sp}:\mathsf P_{\rm CT}^{+}&\longrightarrow
     \mathsf{SurvPath}_{\rm HMT},\\
 F_Q:\mathsf P_{\rm CT}^{+}&\longrightarrow\operatorname{Dom}(\mathcal Q),\\
 F_D:\mathsf P_{\rm CT}^{+}&\longrightarrow
     \operatorname{Dom}(\mathfrak D_{\rm HMT}),\\
 F_C:\mathsf P_{\rm CT}^{+}&\longrightarrow
     \operatorname{Dom}(\rho_C^{\rm disc}).
 \end{aligned}
\]
En particular, \(F_{\rm sp}\) conserva la familia compatible de
restricciones finitas de la sección y de la ruta; no reconstruye esa familia
a partir de su espectro.
\end{definition}

\begin{lemma}[Funtorialidad de los adaptadores]
\label{lem:int68-adaptadores}
Los cuatro adaptadores preservan identidades y composición.
\end{lemma}

\begin{proof}
La identidad de un objeto del producto fibrado es la cuádrupla de identidades
sobre el mismo estado TPK; cada proyección devuelve la identidad del factor.
La composición está definida componente a componente entre cuádruplas que
proyectan sobre la misma concatenación de rutas. Proyectar antes o después de
componer da, por definición, el mismo resultado.
\end{proof}

\subsection{\(4\) componentes y producto positivo}

Los teoremas de realización espectral, hilbertiana, dimensional y de Cartan
proporcionan
\[
 \mathfrak S:\mathsf{SurvPath}_{\rm HMT}
      \longrightarrow\mathbf{SpecRoute},
 \qquad
 \mathcal Q:\operatorname{Dom}(\mathcal Q)
      \longrightarrow\mathsf{Hilb}_{\mathbb R,J},
\]
\[
 \mathfrak D_{\rm HMT}:\operatorname{Dom}(\mathfrak D_{\rm HMT})
      \longrightarrow\mathbf B\mathfrak M_D^+,
 \qquad
 \rho_C^{\rm disc}:\operatorname{Dom}(\rho_C^{\rm disc})
      \longrightarrow\mathbf B G_C,
\]
donde
\[
 G_C:=\operatorname{Spin}^{+}(1,3)\ltimes\mathbb R^{1,3}
\]
y
\begin{equation}
 \mathfrak M_D^+
 :=\mathbb N\times\mathbb R_{\geq0}^{2}\times\mathbb R_{>0},
 \qquad
 (n,a,b,\lambda)\odot(n',a',b',\lambda')
 =(n+n',a+a',b+b',\lambda\lambda').
 \label{eq:int68-monoide-dimensional-positivo}
\end{equation}
Aquí \(\mathbf BM\) es la categoría de un objeto asociada al monoide
\(M\). El primer funtor es el del
\cref{thm:functor-realizacion-espectral-numero}; el segundo es la
realización de Hilbert definida en el capítulo de postulados; el tercero es
el \cref{thm:funtor-dimensional-minimo}; y el cuarto es el
\cref{thm:realizacion-discreta-cartan-rev6}.

Definimos las restricciones comunes
\[
 \mathfrak S_+:=\mathfrak S\circ F_{\rm sp},\qquad
 \mathcal Q_+:=\mathcal Q\circ F_Q,\qquad
 \mathfrak D_+:=\mathfrak D_{\rm HMT}\circ F_D,\qquad
 \rho_{C,+}^{\rm disc}:=\rho_C^{\rm disc}\circ F_C.
\]

\begin{theorem}[Realización discreta positiva sobre el dominio común]
\label{thm:int68-producto-positivo}
La aplicación
\begin{equation}
 \mathfrak R_{\rm MD}^{\rm disc,+}
 :=(\mathfrak S_+,\mathcal Q_+,\mathfrak D_+,
       \rho_{C,+}^{\rm disc})
 \label{eq:int68-producto-positivo}
\end{equation}
es un funtor
\[
 \mathsf P_{\rm CT}^{+}\longrightarrow
 \mathbf{SpecRoute}\times\mathsf{Hilb}_{\mathbb R,J}
 \times\mathbf B\mathfrak M_D^+\times\mathbf BG_C.
\]
No se afirma que toda ruta TPK pertenezca a \(\mathsf P_{\rm CT}^{+}\).
\end{theorem}

\begin{proof}
Por el \cref{lem:int68-adaptadores}, cada proyección \(F_i\) es un funtor;
por los teoremas citados, también lo son
\(\mathfrak S,\mathcal Q,\mathfrak D_{\rm HMT}\) y
\(\rho_C^{\rm disc}\) en sus dominios declarados. Sus composiciones
preservan, pues, identidades y concatenación. La composición en el codominio
producto se calcula componente a componente, lo que prueba
\eqref{eq:int68-producto-positivo}.
\end{proof}

La componente de Hilbert tiene por blanco general
\(\mathsf{Hilb}_{\mathbb R,J}\), cuyos morfismos son operadores reales
acotados que entrelazan las estructuras complejas. Sea
\(\mathsf P_{\rm CT}^{u}\subseteq\mathsf P_{\rm CT}^{+}\) la subcategoría
de las flechas \(\gamma:x\to y\) para las cuales
\begin{equation}
 \mathcal Q_+(\gamma)^*\mathcal Q_+(\gamma)=I_{\mathcal Q_+(x)},
 \qquad
 \mathcal Q_+(\gamma)\mathcal Q_+(\gamma)^*=I_{\mathcal Q_+(y)}.
 \label{eq:int68-unitariedad-doble}
\end{equation}
Las identidades y los productos de operadores unitarios vuelven a ser
unitarios; por ello esta condición define una subcategoría y sólo en ella se
obtiene
\begin{equation}
 \mathcal U_{\rm reg}:=
 \mathcal Q_+\big|_{\mathsf P_{\rm CT}^{u}}:
 \mathsf P_{\rm CT}^{u}\longrightarrow
 \mathsf{UHilb}_{\mathbb R,J}.
 \label{eq:int68-restriccion-unitaria}
\end{equation}
La invertibilidad unilateral o la conservación de norma en un subconjunto no
sustituyen las dos identidades de \eqref{eq:int68-unitariedad-doble}.

La notación de siete entradas se descompone en las cuatro componentes
tipadas: \(Q\) es \(\mathcal Q_+\);
\(U\) es la restricción \eqref{eq:int68-restriccion-unitaria}; \(A\) es la
coordenada de acción de \(\mathfrak D_+\); \(C\) es la conjugación declarada
en el subdominio donde entrelaza la componente hilbertiana y la orientación;
\(\mathcal E\) y \(\mathcal B\) son lectoras constitutivas subordinadas a
los datos dimensionales; y \(N_{\rm sp}\) es una propiedad cuadrática de la
realización partida. Ninguna de las cuatro últimas se añade al producto como
un quinto funtor sin dominio y codominio.

\subsection{Lecturas positiva y orientada}

Sea \(e\mapsto e^\vee\) la reversión de aristas del grafo dirigido
subyacente. En \(\mathsf P_{\rm CT}^{+}\), \(e^\vee\) es otra arista y no
se declara todavía inversa categorial de \(e\). Fijado un peso
\(w(e)=w(e^\vee)>0\), definimos sobre una ruta
\(\gamma=e_1\cdots e_r\)
\[
 V_+(\gamma):=\sum_{j=1}^r w(e_j),
 \qquad
 V_{\rm or}(\gamma):=
 \sum_{j=1}^r\varepsilon(e_j)w(\underline e_j),
\]
donde \(\underline e_j\) es la arista no orientada y
\(\varepsilon(e^\vee)=-\varepsilon(e)\). La primera lectora conserva
variación total positiva; la segunda, transporte neto orientado. Ambas son
aditivas por concatenación y reciben la misma ruta TPK.

\begin{proposition}[No factorización del lector positivo por el orientado]
\label{prop:int68-no-factorizacion}
Supóngase que la identidad y el lazo
\(ee^\vee:=e^\vee\circ e\) pertenecen al dominio construido. No existe una
función \(H\) sobre \(\operatorname{im}V_{\rm or}\) tal que
\(V_+=H\circ V_{\rm or}\).
\end{proposition}

\begin{proof}
Se tiene
\[
 V_{\rm or}(\operatorname{id})=0
 =V_{\rm or}(ee^\vee),
 \qquad
 V_+(\operatorname{id})=0,
 \qquad
 V_+(ee^\vee)=2w(e)>0.
\]
Una misma entrada \(0\) obligaría a \(H(0)\) a tomar dos valores distintos.
\end{proof}

Este resultado es una no factorización concreta en el dominio construido; no
es un teorema universal sobre cualquier entrelazador enriquecido posible.
En particular, no autoriza a borrar una lectora. La unidad está en el
antecedente TPK común y no en una compresión horizontal entre sus salidas.

\subsection{Localización reversible y realización orientada}

Sea \(\mathsf P_{{\rm CT},{\rm rev}}^+\) la subcategoría generada por las
aristas para las cuales se han comprobado las cuatro condiciones siguientes:
\begin{enumerate}
 \item \(\mathfrak S_+(e)\) es un isomorfismo de
       \(\mathbf{SpecRoute}\) y
       \(\mathfrak S_+(e^\vee)=\mathfrak S_+(e)^{-1}\);
 \item \(\mathcal Q_+(e)\) satisface
       \eqref{eq:int68-unitariedad-doble} y
       \(\mathcal Q_+(e^\vee)=\mathcal Q_+(e)^*\);
 \item el dato dimensional orientado procede de una cocadena aditiva con
       \(d_{\rm or}(e^\vee)=-d_{\rm or}(e)\) y toma valores en
       \(\operatorname{Gr}(\mathfrak M_D^+)\);
 \item los cociclos de Cartan son normalizados, aditivos y antisimétricos
       bajo reversión.
\end{enumerate}
Formamos el grupoide
\[
 \Pi_\vee(\mathsf P_{{\rm CT},{\rm rev}}^+)
\]
obtenido de la completación en grupoide imponiendo además
\([e^\vee]=[e]^{-1}\). La componente dimensional orientada toma valores en
\begin{equation}
 \operatorname{Gr}(\mathfrak M_D^+)
 \cong\mathbb Z\times\mathbb R^2\times\mathbb R_{>0},
 \label{eq:int68-grothendieck-dimensional}
\end{equation}
con inversa
\((n,a,b,\lambda)^{-1}=(-n,-a,-b,\lambda^{-1})\). Este lector firmado no es
la magnitud positiva \(V_+\) rebautizada: sus valores sobre las aristas
reversas son distintos.

Para la componente Cartan, sean \(c_g,c_s\) cociclos como en el punto 4,
sea \(R_4\in\operatorname{Spin}^{+}(1,3)\), sea \(R_4u=u\) y sea
\(\ell_0>0\). Entonces
\[
 \rho_C^{\rm disc}(\gamma)
 =\bigl(R_4^{c_g(\gamma)},\ell_0c_s(\gamma)u\bigr).
\]
La ley
\((R,a)(S,b)=(RS,a+Rb)\), la igualdad \(R_4u=u\) y la aditividad de los
cociclos dan
\[
 \rho_C^{\rm disc}(\gamma_2\gamma_1)
 =\rho_C^{\rm disc}(\gamma_2)\rho_C^{\rm disc}(\gamma_1),
 \qquad
 \rho_C^{\rm disc}(\gamma^\vee)
 =\rho_C^{\rm disc}(\gamma)^{-1}.
\]

\begin{theorem}[Realización discreta orientada]
\label{thm:int68-realizacion-orientada}
Bajo las cuatro condiciones anteriores, las componentes descienden de manera
única a un funtor
\[
 \mathfrak R_{\rm MD}^{\rm disc,or}:
 \Pi_\vee(\mathsf P_{{\rm CT},{\rm rev}}^+)
 \longrightarrow
 \operatorname{Core}(\mathbf{SpecRoute})
 \times\mathsf{UHilb}_{\mathbb R,J}
 \times\mathbf B\operatorname{Gr}(\mathfrak M_D^+)
 \times\mathbf BG_C.
\]
En particular,
\[
 \mathfrak R_{\rm MD}^{\rm disc,or}(\gamma^{-1})
 =\mathfrak R_{\rm MD}^{\rm disc,or}(\gamma)^{-1}.
\]
\end{theorem}

\begin{proof}
Las condiciones 1 y 2 hacen invertibles las imágenes espectral e
hilbertiana; la condición 3 toma valores en un grupo; y la componente Cartan
ya toma valores en \(G_C\). Las cuatro respetan la relación generadora
\([e^\vee]=[e]^{-1}\). Por la propiedad universal de la localización en
grupoide, cada una desciende de manera única; su producto es el funtor
indicado. La identidad de inversión se verifica componente a componente.
\end{proof}

La lectora positiva \(\mathfrak R_{\rm MD}^{\rm disc,+}\) y la orientada
\(\mathfrak R_{\rm MD}^{\rm disc,or}\) son hermanas: ninguna se define como
compresión de la otra. El paso posterior de un valor orientado a su módulo no
es, en general, monoidal y coincide con una lectura positiva sólo en un cono
sin cancelación. Del mismo modo, este par no sustituye la doble proyección
fundacional: la evaluación arquimediana y la rama
Paley--Witt--Golay--Leech--\(V^\natural\)--Monstruo--Moonshine permanecen
como lectoras hermanas del mismo estado enriquecido. Las cuatro componentes
construidas en esta sección no sustituyen ninguna de esas dos proyecciones.

\subsection{Entrelazamiento de Cartan y límite suave}

Sea
\[
 \mathcal A_{\ell_0}
 =\begin{pmatrix}\omega&\ell_0^{-1}e\\0&0\end{pmatrix}
\]
una conexión de Cartan declarada. La ecuación
\begin{equation}
 \operatorname{Hol}_{\mathcal A_{\ell_0}}
   \bigl(\mathfrak R_C(\gamma)\bigr)
 =\rho_C\bigl(\operatorname{Hol}_{\rm TPK}(\gamma)\bigr)
 \label{eq:int68-cuadrado-cartan}
\end{equation}
es exacta en el modelo discreto cuando ambas holonomías se definen por el
mismo producto ordenado de aristas.

\begin{criterion}[Promoción a una realización suave]
\label{crit:int68-cartan-suave}
La ecuación \eqref{eq:int68-cuadrado-cartan} sólo define una prolongación
suave si se prueban, además, independencia de subdivisión, convergencia del
producto ordenado y recuperación local de una cotétrada y una conexión con
los dominios y regularidad declarados.
\end{criterion}

El término «físico» designa aquí que las cuatro componentes publican
espectro, realización de estado, dimensiones y transporte de Cartan sin
perder su antecedente TPK\@. No afirma unicidad, equivalencia de categorías,
cobertura de toda ruta, carta SI automática ni teoría de campos completa.

\begin{criterion}[Obstrucciones a las realizaciones positiva, unitaria y orientada]
\label{crit:int68-obstrucciones-realizacion}
La realización positiva no está definida si falta alguno de los cuatro
adaptadores, si sus imágenes no poseen la misma ruta TPK subyacente o si una
componente no preserva identidad y composición. La restricción unitaria no
existe si falta cualquiera de las dos igualdades de
\eqref{eq:int68-unitariedad-doble}. La realización orientada no desciende al
grupoide si una imagen espectral no es isomorfa, si una relación localizada
cambia de imagen, si una magnitud positiva se trata como inversa, si una
reversión no cambia el signo de los datos orientados o si el cociclo de
Cartan deja de ser normalizado, aditivo o antisimétrico. No existe
prolongación suave compatible si el resultado depende de la subdivisión o si
el límite no recupera cotétrada y conexión. Cualquier uso de una coordenada
externa para seleccionar la ruta viola el dominio construido.
\end{criterion}

\paragraph{Dominio, codominio y estatuto de las ecuaciones de confluencia HMT–MD}
El producto positivo queda demostrado sobre el producto fibrado común y la
realización orientada, sobre la localización reversible bajo las cuatro
condiciones explícitas. Extender el dominio a otras rutas exige construir sus
cuatro levantamientos sobre un mismo antecedente. La prolongación suave exige,
además, las propiedades de subdivisión, convergencia y recuperación local del
\cref{crit:int68-cartan-suave}; una verificación finita no sustituye esas
propiedades.

\endgroup

\paragraph{Paso de la realización Cartan–Holst a la confluencia horizonte–área–información mediante conservación de la acción} La primera realización constitutiva de esta confluencia es el vacío electromagnético--electrodébil, donde las funciones simétrica, antisimétrica y determinantal publican propagación y acoplamientos.
